% GENERATED FILE. Edit canonical lessons, then run npm run generate:books.
% canonical-source-hash: fnv1a64:f7e541c1d3a77399
% canonical-lessons: SA-C352-himsa, SA-C352-sahisnuh, SA-C352-apakaroti, SA-C352-anugrhnati, SA-C352-ksama

\chapter{Harm and No Harm}
\label{ch:sa-harm-and-no-harm}
\hlchaptermodality{352}

\begin{chapteropening}
\textbf{By the end of this chapter:} \emph{I can talk about violence, harm and forbearance.}

Complete the last lesson of chapter 352: I can talk about violence, harm and forbearance.
\end{chapteropening}

\section[\texorpdfstring{\sk{हिंसा} (hiṃsā)}{हिंसा (hiṃsā)}]{\sk{हिंसा} (hiṃsā) — violence; harm}

\label{lesson:SA-C352-himsa}



\begin{quote}
Before the new one: say the Sanskrit for abides by, then the Sanskrit for the bounds of propriety.
\end{quote}

\subsection*{You'll want to know: \sk{हिंसा}}
\textbf{\sk{हिंसा}} — \emph{hiṃsā} — \textquotedblleft{}violence; harm\textquotedblright{}.

\begin{quote}
\sk{हिंसा} \sk{कदापि} \sk{उचिता} \sk{न} \sk{भवति।}
\end{quote}

\begin{quote}
\emph{hiṃsā kadāpi ucitā na bhavati.}
\end{quote}

\begin{quote}
Violence is not right at any time.
\end{quote}

A feminine noun.

\begin{cousinweb}
\textbf{\sk{हिंसा}} (\emph{hiṃsā}) comes from the root √hiṃs, \textquotedblleft{}to injure\textquotedblright{}; MW gives \emph{hiṃsā} \textquotedblleft{}injury, harm, hurt, mischief\textquotedblright{} (MW).
\end{cousinweb}

\begin{grammarlens}[title={What you've built}]
One more word for giving reasons and opinions.
\end{grammarlens}

\subsection*{Your turn}
\begin{itemize}
\raggedright
  \item \emph{Say it:} \emph{hiṃsā}
  \item \emph{Say it:} \emph{hiṃsā}, once more
  \item \emph{Say it:} \emph{hiṃsā kadāpi ucitā na bhavati.} — Violence is not right at any time.
  \item \emph{Recall:} say the Sanskrit for abides by, then the Sanskrit for the bounds of propriety, then say \emph{hiṃsā} again
\end{itemize}

\begin{tcolorbox}[breakable,colback=teal!4,colframe=teal!35!black,title={Before you move on}]
What does \sk{हिंसा} mean? (\textquotedblleft{}Violence; harm\textquotedblright{}.) Say it once more.
\end{tcolorbox}
\section[\texorpdfstring{\sk{सहिष्णुः} (sahiṣṇuḥ)}{सहिष्णुः (sahiṣṇuḥ)}]{\sk{सहिष्णुः} (sahiṣṇuḥ) — tolerant; forbearing}

\label{lesson:SA-C352-sahisnuh}



\begin{quote}
Before the new one: say the Sanskrit for the bounds of propriety, then the Sanskrit for violence.
\end{quote}

\subsection*{You'll want to know: \sk{सहिष्णुः}}
\textbf{\sk{सहिष्णुः}} — \emph{sahiṣṇuḥ} — \textquotedblleft{}tolerant; forbearing\textquotedblright{}.

\begin{quote}
\sk{सहिष्णुः} \sk{जनः} \sk{क्रोधं} \sk{न} \sk{करोति।}
\end{quote}

\begin{quote}
\emph{sahiṣṇuḥ janaḥ krodhaṃ na karoti.}
\end{quote}

\begin{quote}
A tolerant person does not give way to anger.
\end{quote}

\begin{cousinweb}
\textbf{\sk{सहिष्णुः}} (\emph{sahiṣṇuḥ}) comes from the root √sah, \textquotedblleft{}to bear, to put up with\textquotedblright{}, with the ending \emph{-iṣṇu}; MW gives \emph{sahiṣṇu} \textquotedblleft{}patient, forbearing\textquotedblright{} and \textquotedblleft{}bearing, enduring\textquotedblright{}, and KRM lists \emph{sahiṣṇu} among the forms of the root (MW).
\end{cousinweb}

\begin{grammarlens}[title={What you've built}]
One more word for giving reasons and opinions.
\end{grammarlens}

\subsection*{Your turn}
\begin{itemize}
\raggedright
  \item \emph{Say it:} \emph{sahiṣṇuḥ}
  \item \emph{Say it:} \emph{sahiṣṇuḥ}, once more
  \item \emph{Say it:} \emph{sahiṣṇuḥ janaḥ krodhaṃ na karoti.} — A tolerant person does not give way to anger.
  \item \emph{Recall:} say the Sanskrit for the bounds of propriety, then the Sanskrit for violence, then say \emph{sahiṣṇuḥ} again
\end{itemize}

\begin{tcolorbox}[breakable,colback=teal!4,colframe=teal!35!black,title={Before you move on}]
What does \sk{सहिष्णुः} mean? (\textquotedblleft{}Tolerant; forbearing\textquotedblright{}.) Say it once more.
\end{tcolorbox}
\section[\texorpdfstring{\sk{अपकरोति} (apakaroti)}{अपकरोति (apakaroti)}]{\sk{अपकरोति} (apakaroti) — does harm to; wrongs}

\label{lesson:SA-C352-apakaroti}



\begin{quote}
Before the new one: say the Sanskrit for violence, then the Sanskrit for tolerant.
\end{quote}

\subsection*{You'll want to know: \sk{अपकरोति}}
\textbf{\sk{अपकरोति}} — \emph{apakaroti} — \textquotedblleft{}does harm to; wrongs\textquotedblright{}.

\begin{quote}
\sk{सः} \sk{कस्य} \sk{अपि} \sk{न} \sk{अपकरोति।}
\end{quote}

\begin{quote}
\emph{saḥ kasya api na apakaroti.}
\end{quote}

\begin{quote}
He does no one any harm.
\end{quote}

Past (he …): \sk{सः} \sk{मम} \sk{अपकृतवान्}, \emph{saḥ mama apakṛtavān}.

\begin{cousinweb}
\textbf{\sk{अपकरोति}} (\emph{apakaroti}) is built on \emph{apa}, \textquotedblleft{}away, off\textquotedblright{}, and the root √kṛ, \textquotedblleft{}to do\textquotedblright{}; MW gives \emph{apa-kṛta} \textquotedblleft{}done wrongly or maliciously\textquotedblright{} (MW).
\end{cousinweb}

\begin{grammarlens}[title={What you've built}]
One more word for giving reasons and opinions.
\end{grammarlens}

\subsection*{Your turn}
\begin{itemize}
\raggedright
  \item \emph{Say it:} \emph{apakaroti}
  \item \emph{Say it:} \emph{apakaroti}, once more
  \item \emph{Say it:} \emph{saḥ kasya api na apakaroti.} — He does no one any harm.
  \item \emph{Recall:} say the Sanskrit for violence, then the Sanskrit for tolerant, then say \emph{apakaroti} again
\end{itemize}

\begin{tcolorbox}[breakable,colback=teal!4,colframe=teal!35!black,title={Before you move on}]
What does \sk{अपकरोति} mean? (\textquotedblleft{}Does harm to; wrongs\textquotedblright{}.) Say it once more.
\end{tcolorbox}
\section[\texorpdfstring{\sk{अनुगृह्णाति} (anugṛhṇāti)}{अनुगृह्णाति (anugṛhṇāti)}]{\sk{अनुगृह्णाति} (anugṛhṇāti) — shows favour to; obliges}

\label{lesson:SA-C352-anugrhnati}



\begin{quote}
Before the new one: say the Sanskrit for tolerant, then the Sanskrit for does harm to.
\end{quote}

\subsection*{You'll want to know: \sk{अनुगृह्णाति}}
\textbf{\sk{अनुगृह्णाति}} — \emph{anugṛhṇāti} — \textquotedblleft{}shows favour to; obliges\textquotedblright{}.

\begin{quote}
\sk{गुरुः} \sk{शिष्यम्} \sk{अनुगृह्णाति।}
\end{quote}

\begin{quote}
\emph{guruḥ śiṣyam anugṛhṇāti.}
\end{quote}

\begin{quote}
The teacher shows his pupil favour.
\end{quote}

Past (he …): \sk{सः} \sk{माम्} \sk{अनुगृहीतवान्}, \emph{saḥ mām anugṛhītavān}.

\begin{cousinweb}
\textbf{\sk{अनुगृह्णाति}} (\emph{anugṛhṇāti}) is built on \emph{anu}, \textquotedblleft{}after\textquotedblright{}, and the root √grah, \textquotedblleft{}to seize, to take\textquotedblright{}; MW gives \emph{anu-√grah} \textquotedblleft{}to support\textquotedblright{} and \textquotedblleft{}to treat with kindness, to favour, to oblige\textquotedblright{} (MW).
\end{cousinweb}

\begin{grammarlens}[title={What you've built}]
One more word for giving reasons and opinions.
\end{grammarlens}

\subsection*{Your turn}
\begin{itemize}
\raggedright
  \item \emph{Say it:} \emph{anugṛhṇāti}
  \item \emph{Say it:} \emph{anugṛhṇāti}, once more
  \item \emph{Say it:} \emph{guruḥ śiṣyam anugṛhṇāti.} — The teacher shows his pupil favour.
  \item \emph{Recall:} say the Sanskrit for tolerant, then the Sanskrit for does harm to, then say \emph{anugṛhṇāti} again
\end{itemize}

\begin{tcolorbox}[breakable,colback=teal!4,colframe=teal!35!black,title={Before you move on}]
What does \sk{अनुगृह्णाति} mean? (\textquotedblleft{}Shows favour to; obliges\textquotedblright{}.) Say it once more.
\end{tcolorbox}
\section[\texorpdfstring{\sk{क्षमा} (kṣamā)}{क्षमा (kṣamā)}]{\sk{क्षमा} (kṣamā) — forbearance; patience}

\label{lesson:SA-C352-ksama}



\begin{quote}
Before the new one: say the Sanskrit for does harm to, then the Sanskrit for shows favour to.
\end{quote}

\subsection*{You'll want to know: \sk{क्षमा}}
\textbf{\sk{क्षमा}} — \emph{kṣamā} — \textquotedblleft{}forbearance; patience\textquotedblright{}.

\begin{quote}
\sk{क्रोधस्य} \sk{औषधं} \sk{क्षमा} \sk{अस्ति।}
\end{quote}

\begin{quote}
\emph{krodhasya auṣadhaṃ kṣamā asti.}
\end{quote}

\begin{quote}
The medicine for anger is forbearance.
\end{quote}

A feminine noun.

\begin{cousinweb}
\textbf{\sk{क्षमा}} (\emph{kṣamā}) comes from the root √kṣam, \textquotedblleft{}to be patient, to endure\textquotedblright{}; MW gives \emph{kṣamā} \textquotedblleft{}patience, forbearance, indulgence\textquotedblright{} (MW).
\end{cousinweb}

\begin{grammarlens}[title={What you've built}]
One more word for giving reasons and opinions.
\end{grammarlens}

\subsection*{Your turn}
\begin{itemize}
\raggedright
  \item \emph{Say it:} \emph{kṣamā}
  \item \emph{Say it:} \emph{kṣamā}, once more
  \item \emph{Say it:} \emph{krodhasya auṣadhaṃ kṣamā asti.} — The medicine for anger is forbearance.
  \item \emph{Recall:} say the Sanskrit for does harm to, then the Sanskrit for shows favour to, then say \emph{kṣamā} again
\end{itemize}

\begin{tcolorbox}[breakable,colback=teal!4,colframe=teal!35!black,title={Before you move on}]
What does \sk{क्षमा} mean? (\textquotedblleft{}Forbearance; patience\textquotedblright{}.) Say it once more.
\end{tcolorbox}
